\begin{lemma}[Reserved colors at a hierarchy stage]
Let a hierarchy be good.  For every $0\leq j\leq q$ and every
$\eta\in D_{n,j}$,
\[
 \Gamma^j_\eta\subseteq
 \overline{\varphi_n}(T_{n,j}^{*}).
\]
Consequently
\[
 \Gamma^j\cap D_{n,j}=\varnothing,\qquad
 \Gamma^j=\bigcup_{\eta\in D_{n,j}}\Gamma^j_\eta.
\]
\end{lemma}
\begin{proof}
Fix $j$ and $\eta\in D_{n,j}$.  Clause (i), namely the avoidance field
of \noderef{GoodHierarchy}, puts $\Gamma^j_\eta$ in the relevant
stage-missing set of \noderef{HierarchyStageMissingColors}.  For
$j>0$ that tree is exactly $T_{n,j}^{*}=T_{n,j}$.

For $j=0$, clause (i) gives
$\Gamma^0_\eta\subseteq\overline{\varphi_n}(T_n)$.  If the terminal
kind is not PE, \noderef{HierarchyBase} identifies
$T_{n,0}^{*}$ with $T_n$.  If it is PE, the same definition supplies
an \noderef{OrderedTreeUnion} witness whose left-segment field embeds
$T_n$ into $T_n\vee R_n=T_{n,0}^{*}$.  In either case every vertex of
$T_n$ belongs to the starred base.  The union definition
\noderef{TreeMissingColors} therefore gives
\[
 \overline{\varphi_n}(T_n)
 \subseteq\overline{\varphi_n}(T_{n,0}^{*}),
\]
proving the first assertion also at $j=0$.

By \noderef{HierarchyStoredColorsAt},
$D_{n,j}=D_n-\overline{\varphi_n}(T_{n,j}^{*})$.  Thus no color in
$D_{n,j}$ belongs to the starred-stage missing set.  Every member of
every $\Gamma^j_\eta$ does belong to that missing set by the first
assertion, so it cannot lie in $D_{n,j}$.  Taking the union over
$\eta\in D_{n,j}$ gives
$\Gamma^j\cap D_{n,j}=\varnothing$.
\end{proof}
